% Configurazione
\documentclass{article}

\usepackage{titling} % Required for inserting the subtitle
\usepackage{graphicx} % Required for inserting images
\usepackage{tabularx} % Per l'ambiente tabularx (tabelle)
\usepackage{calc} % Sempre per le tabelle
\usepackage[hidelinks]{hyperref} % Per i collegamenti ipertestuali, ad esempio sulla table of contents
\usepackage[italian]{babel} % Per la lingua italiana nelle scritte automatiche
\usepackage{xcolor} % Per colorare il testo
\usepackage{colortbl} % Per colorare le celle delle tabelle
\usepackage{lipsum} % Per generare lorem ipsum
\usepackage[normalem]{ulem} % Per sottolineare il testo
\usepackage{array} % Per la visualizzazione fluttuante di array di domande e risposte
\usepackage{ragged2e} % Pacchetto necessario per \justifying che giustifica il testo di tabelle
\usepackage{tikz} % Per spostare elementi nel documento in modo facile e veloce
\usepackage[a4paper, top=2.5cm, bottom=2.5cm, left=2.5cm, right=2.5cm]{geometry} % Per i margini della pagina
\usepackage{fancyhdr} % Per l'intestazione e il piè di pagina
\usepackage{amsmath} % Per scrivere formule matematiche, in particolare per il pedice G

\newcommand{\ulhref}[2]{\href{#1}{\uline{#2}}} % Nuovo comando per sottolineare i link
\newcommand{\ulref}[1]{\uline{\ref{#1}}} % Nuovo comando per sottolineare i collegamenti a immagini e tabelle
\setlength{\parindent}{0pt} % Rimuove il rientro automatico dei paragrafi
\usetikzlibrary{calc} % Libreria per il calcolo delle coordinate di TikZ
\pagestyle{fancy} % Stile della pagina, per l'intestazione e il piè di pagina
\renewcommand{\footrulewidth}{0.4pt} % Inserimento della linea orizzontale in basso
\setlength{\headsep}{1.4cm} % Spazio tra l'intestazione e il testo

\graphicspath{ {immagini/} {../../../../shared/immagini/} }



% Struttura
\begin{document}

% Parte costante
\thispagestyle{plain} % Niente intestazione e piè di pagina


\begin{tikzpicture}[remember picture, overlay]
    % Punto di partenza al centro orizzontale nella metà superiore
    \coordinate (top_center) at ($(current page.north)!0.3!(current page.south)$);

    % UniPD: Logo e descrizione
    \node at (top_center) [anchor=north, xshift=-3cm, yshift=4.85cm] 
        {\includegraphics[width=0.15\textwidth]{Logo Universita di Padova.png}};
    \node at (top_center) [anchor=north, xshift=1.7cm, yshift=4.5cm]
        {\textcolor{red}{\textbf{Università degli Studi di Padova}}};
    \node at (top_center) [anchor=north, xshift=1.7cm, yshift=4cm]
        {\textcolor{red}{Laurea: Informatica}};
    \node at (top_center) [anchor=north, xshift=1.7cm, yshift=3.5cm]
        {\textcolor{red}{Corso: Ingegneria del Software}};
    \node at (top_center) [anchor=north, xshift=1.7cm, yshift=3cm]
        {\textcolor{red}{Anno Accademico: 2024/2025}};

    % SWEg Labs: Logo e descrizione
    \node at (top_center) [anchor=north, xshift=-2.85cm, yshift=1.5cm] 
        {\includegraphics[width=0.16\textwidth]{Logo SWEg.png}};
    \node at (top_center) [anchor=north, xshift=1.7cm, yshift=0.5cm]
        {\textbf{Gruppo: SWEg Labs}};
    \node at (top_center) [anchor=north, xshift=1.7cm, yshift=0cm]
        {Email: \textsf{gruppo.sweg@gmail.com}};
\end{tikzpicture}


\vspace{10cm}

{
\centering
\Huge\bfseries Norme di progetto\par
\vspace{0.5cm}
\Large Versione 1.0.0\par
}

\vspace{2cm}
\newpage
% Intestazione
\fancyhead[L]{Indice} % Testo a sinistra

\tableofcontents
\section{Introduzione}
\subsection{Scopo del documento}
\subsection{Scopo del prodotto}
\subsection{Glossario}
\subsection{Maturità e miglioramenti}
\subsection{Riferimenti}
\subsubsection{Riferimenti normativi}
\subsubsection{Riferimenti informativi}

\section{Analisi dei rischi}
\subsection{RT: Rischi legati alle tecnologie}
\subsubsection{Complessità delle nuove tecnologie}
\subsubsection{Mancanza di risorse e documentazione}
\subsubsection{Aggiornamenti o modifiche agli strumenti e tecnologie in uso}
\subsection{RO: Rischi legati all'organizzazione del gruppo}
\subsubsection{Rischi di comunicazione interna}
\subsubsection{Rischi di confusione sulle responsabilità}
\subsubsection{Rischi legati alla gestione del tempo e delle scadenze}
\subsection{RP: Rischi legati ai singoli membri del gruppo}
\subsubsection{Rischi legati alla continuità del progetto}
\subsubsection{Rischi legati alla non conformità rispetto agli impegni dichiarati}
\newpage

% Parte variabile
\begin{center}
\begin{tabular}{r|l}
    \textbf{Versione} & 1.0 \\
    \hline
    \textbf{Stato} & Approvato \\
    \hline
    \textbf{Uso} & Esterno \\
    \hline
    \textbf{Approvazione} & \emph{Tutto il gruppo} \\
    \hline
    \textbf{Redazione} & Federica Bolognini \\
    \hline
    \textbf{Verifica} & ....... \\
    \hline
    \textbf{Distribuzione} & \emph{SWEg Labs} \\
    \hline
    \textbf{Destinatari} & Prof. Tullio Vardanega \\
                         & Prof. Riccardo Cardin \\
\end{tabular}
\end{center}

\newpage
\section{Ordine del giorno}

\begin{itemize}
    \item Preparare le domande per l'azienda \emph{Vimar} in un Foglio Google, come da loro richiesto, cosicchè possano rispondere a distanza.
    \item Creazione di un template di presentazione per il diario di bordo, il cui primo incontro è previsto per \emph{lunedì 28-10-24}.
    \item Organizzazione del gruppo per quanto riguarda la documentazione che è necessario consegnare alla candidatura (cioè entro \emph{venerdì 01-11-24} alle ore 17):
    \begin{itemize}
        \renewcommand{\labelitemii}{--}
        \item Lettera di presentazione
        \item Preventivo degli impegni
        \item Valutazione dei capitolati
        \item Verbali delle riunioni interne ed esterne
    \end{itemize}
\end{itemize}
\newpage
\section{Diario della riunione}

\begin{itemize}
    \item Discussione con il proponente \emph{Ergon} sull'artefatto richiesto, guidata dalle nostre domande.
    \item Risposta del proponente ad alcune domande:
\end{itemize}
\vspace{1cm}
\begingroup
\renewcommand{\ni}{\noindent}

{\Large Domande funzionali}

\vspace{0.5cm}

\begin{tabular}{>{\justifying\arraybackslash}p{0.4\textwidth} >{\justifying\arraybackslash}p{0.6\textwidth}}
    \multicolumn{1}{c}{\textbf{Domande}} & \multicolumn{1}{c}{\textbf{Risposte}} \\ \\
    \ni Bisogna costruire il modello di embedding da zero oppure è possibile sfruttare langchain? Nel caso quali tecnologie sono consigliate? & \ni Si possono usare anche cose già fatte, integrandole nel progetto. Libertà piena \\ \\
    \ni Cosa intendete per template? nello specifico la configurazione dei template delle domande nel back-end deve essere aperta alla modifica in caso diventassero antiquati? & \ni Ci sono due possibili tipologie di domande: template di domanda e risposta, e la totale IA generativa. Se ci sono feedback, l'obbiettivo è andare mano mano a creare dei template di domande frequenti a risposta precisa. \\ \\
    \ni Come dovrà essere fatta l'interfaccia utente per la configurazione della piattaforma? & \ni Un'unica chat. Interessante la cronologia, ma non essenziale. \\ \\
    \ni Come viene raccolto e utilizzato il feedback dell'utente per migliorare le risposte future nel sistema? & \ni Bisogna valutare in base all'LLM preaddestrato da cui attingeremo. In fase iniziale, va bene anche solo il pollice in su e in giù. In fase avanzata, il chatbot può anche chiedere all'utente che cosa si aspettava \\ \\
    \ni Come deve essere visualizzato il prompt e strutturato l’output? deve essere anche desktop o solo mobile? & \ni  Idea era come un'app mobile, preferenza per .NET MAUOI. Sono aperti anche ad una web app. Focus comunque sul mobile \\ \\
\end{tabular}
\begin{tabular}{>{\justifying\arraybackslash}p{0.4\textwidth} >{\justifying\arraybackslash}p{0.6\textwidth}}
    \ni Quali sono le funzionalità richieste lato front-end? deve essere solamente chat testuale o deve permettere il caricamento di un file? & \ni Semplicemente chat \\ \\
    \ni Come deve essere gestito il login? dovrà avere un'utenza singola o multipla? & \ni Un account per membro di azienda, quindi multipli \\ \\
    \ni È richiesto di memorizzare le chat con il bot? se si, quante interazioni? & \ni Non è obbligatorio salvare uno storico \\ \\

\end{tabular}

\vspace{1.5cm}

{\Large Domande non funzionali}

\vspace{0.5cm}

\begin{tabular}{>{\justifying\arraybackslash}p{0.4\textwidth} >{\justifying\arraybackslash}p{0.6\textwidth}}
    \multicolumn{1}{c}{\textbf{Domande}} & \multicolumn{1}{c}{\textbf{Risposte}} \\ \\
    \ni Bisogna limitare gli argomenti che si possono trattare con il sistema? & \ni Loro forniranno un contesto. Ideale sarebbe limitarlo lì, non lasciarlo troppo libero \\ \\
    \ni Ci potreste descrivere i dati aziendali che dobbiamo utilizzare ed il modo a cui accedervi? & \ni Sì, verrà fornito un accesso in futuro \\ \\
    \ni Ci potete fornire dei comportamenti tipo in modo da istruire il bot?(nello specifico, definiamo delle regole che il bot deve seguire) & \ni Sì, probabilmente mondo del beverage, aziende che distrubuiscono bevande \\ \\
    \ni Ci potreste dare APIkey per i vostri strumenti? & \ni Forniranno aiuto anche con strumenti aziendali. Mentre i dati vengono forniti tramite CSV \\ \\
\end{tabular}

\vspace{1.5cm}

{\Large Altro}
\nopagebreak
\vspace{0.5cm}

\begin{tabular}{>{\justifying\arraybackslash}p{0.4\textwidth} >{\justifying\arraybackslash}p{0.6\textwidth}}
    \multicolumn{1}{c}{\textbf{Domande}} & \multicolumn{1}{c}{\textbf{Risposte}} \\ \\
    \ni Quali sono le tempistiche per i corsi online che ci mettete a disposizione? & \ni I corsi sono offline, li possiamo consultare quando vogliamo. Ce ne sarà uno su .NET MAUI e anche su LLM. Tutta documentazione offline \\ \\
    \ni Possiamo utilizzare qualsiasi LLM Esterno? & \ni Possiamo scegliere quello che vogliono. Anche non open source, ma lo dobbiamo chiedere. \\ \\
\end{tabular}


\endgroup
\newpage
% Intestazione
\fancyhead[L]{4 \hspace{0.2cm} Decisioni} % Testo a sinistra


\section{Decisioni}

Durante la riunione sono state prese le seguenti decisioni:

\vspace{0.5cm}

\begin{table}[htbp]
    \centering
    \rowcolors{2}{lightgray}{white}
    \begin{tabular}{|c|p{0.8\textwidth}|}
        \hline
        \rowcolor[gray]{0.75}
        \multicolumn{1}{|c|}{\textbf{Codice}} & \multicolumn{1}{|c|}{\textbf{Descrizione}}\\
        \hline
        VI 8.1 & È stato deciso di incominciare a scrivere i documenti per la RTB, in alcune parti fattibili e prestabilite, entro lunedì 11 novembre,
        dove svolgeremo una riunione interna nella quale prepareremo l'incontro con \emph{AzzurroDigitale} del giorno dopo \\
        \hline
        VI 8.2 & Sono state decise le metriche di qualità riguardanti la documentazione, cioè gli obiettivi di correttezza linguistica e leggibilità.
        In particolare, abbiamo stabilito per gli errori ortografici un valore accettabile di 5\% e un valore ottimale di 0\%, e per la leggibilità abbiamo
        stabilito di utilizzare l'Indice di Gulpease, con un valore accettabile di 40 e un valore ottimale di 60 \\
        \hline
        VI 8.3 & È stato deciso di utilizzare \emph{GitHub Projects} per la gestione del backlog, poichè esso è già incluso nel nostro ambiente di lavoro e ci
        permette di avere un'interfaccia unica per tutto il progetto \\
        \hline
        VI 8.4 & È stato deciso di utilizzare \emph{Fogli Google} per creare tutti i grafici e alcuni diagrammi legati al progetto, poichè lo strumento 
        facilita l'inserimento e la gestione di grandi quantità di dati, ed è inoltre già mediamente conosciuto dai membri del gruppo, a differenza ad 
        esempio di uno strumento come \emph{Jira} \\
        \hline
        VI 8.5 & È stato deciso di mantenere due repo separati, uno per la documentazione e uno per il codice, e di utilizzare \emph{GitHub Projects} per la
        gestione del backlog di entrambi i repo \\
        \hline
    \end{tabular}
    \end{table}

\newpage
\section{Todo}

Durante la riunione sono emersi i seguenti task da svolgere:

\vspace{0.5cm}

\begin{table}[htbp]
\begin{tabular}{|p{0.4\textwidth}|p{0.6\textwidth}|}
    \hline
    \rowcolor[gray]{0.9}
    \multicolumn{1}{|c|}{\textbf{Assegnatario}} & \multicolumn{1}{|c|}{\textbf{Task Todo}} \\
    \hline
    \emph{[Tutto il gruppo, singolarmente]} & Approfondimento competenze \emph{LaTeX} \\
    \hline
    \emph{[Tutto il gruppo, singolarmente]} & Approfondimento competenze \emph{Git}, \emph{GitHub} \\
    \hline
    Stefani Riccardo & Redigere verbale riunione \emph{16/10/24} \\
    \hline
\end{tabular}
\end{table}


\end{document}
